% Industry résumé (1 page). Build: pdflatex Jose_Carter_Resume_Industry_2026.tex
\documentclass[letterpaper,10pt]{article}
% Shared style for José Carter's résumés (matches the 2026 WDAS résumé look).
\usepackage[utf8]{inputenc}
\usepackage[T1]{fontenc}
\usepackage{lmodern}
\usepackage[english]{babel}
\usepackage[letterpaper,top=0.45in,bottom=0.45in,left=0.5in,right=0.5in]{geometry}
\usepackage{titlesec}
\usepackage[usenames,dvipsnames]{xcolor}
\usepackage{enumitem}
\usepackage{tabularx}
\usepackage[hidelinks,pdfusetitle]{hyperref}
\usepackage{microtype}

\definecolor{cvblue}{HTML}{1F3F73}
\pagestyle{empty}
\setlength{\parindent}{0pt}
\setlength{\tabcolsep}{0in}
\raggedright
\urlstyle{same}

\titleformat{\section}{\color{cvblue}\Large\bfseries}{}{0em}{}[\vspace{-2pt}\color{cvblue}\titlerule]
\titlespacing*{\section}{0pt}{4pt}{3pt}

\newcommand{\cvheader}[2]{%
  \begin{center}
    {\color{cvblue}\fontsize{22}{24}\selectfont\bfseries José Carter Arriagada}\\[2pt]
    {\small\bfseries #1}\\[2pt]
    {\small Santiago, Chile \textbar{} Chile: +56 9 8655 2408 \textbar{} U.S.: +1 276 221 3064}\\
    {\small \href{mailto:jose.carterx@gmail.com}{jose.carterx@gmail.com} \textbar{} \href{https://josecarter.dev}{josecarter.dev} #2}\\
    {\small \href{https://linkedin.com/in/jose-carter-arriagada}{linkedin.com/in/jose-carter-arriagada} \textbar{} \href{https://github.com/Cartterr}{github.com/Cartterr}}
  \end{center}
  \vspace{-10pt}
}

% \role{Title}{Dates}{Organization}{Location}
\newcommand{\role}[4]{%
  \vspace{1pt}
  \begin{tabular*}{\textwidth}{l@{\extracolsep{\fill}}r}
    \textbf{#1} & \small #2 \\
    \textit{\small #3} & \textit{\small #4} \\
  \end{tabular*}\vspace{-1pt}
}
\newenvironment{points}{\begin{itemize}[leftmargin=0.18in,topsep=0pt,itemsep=0pt,parsep=0pt,partopsep=0pt,label={\small\textbullet}]\small}{\end{itemize}\vspace{1pt}}
\newcommand{\entry}[2]{{\small\textbf{#1} --- #2}\par\vspace{1pt}}

\hypersetup{pdftitle={José Carter Arriagada - Industry Resume 2026}}

\begin{document}
\cvheader{Software Engineer \textbar{} Real-Time XR, Physical AI \& GPU Computing}{}

\section{Profile}
{\small Software engineer (PUC Chile, 2025) building performance-critical software: real-time XR and AR research systems at DICTUC (full time) and physical-AI robotics software at TEK Chile (part-time freelance). Track record of measured speedups on GPUs and parallel systems: 15x in CUDA simulation and 10x in data pipelines.\par}

\section{Technical Skills}
{\small
\textbf{Languages:} Python, TypeScript/JavaScript, C++, C\#, Swift, Java, SQL \quad \textbf{English:} Advanced (TOEIC 905)\\
\textbf{GPU \& Real-Time:} CUDA, Metal, Unity, OpenXR, ARKit/RealityKit, eye and hand tracking, Isaac Sim, performance profiling\\
\textbf{Robotics \& Vision:} YOLO, PyTorch, robot-arm and mobile-robot control, Modbus, TCP/IP, Raspberry Pi\\
\textbf{Platforms:} Node.js, React, Angular/Nx, Vue, Next.js, PostgreSQL, InfluxDB, MQTT, AWS, Docker, Railway, Linux, CI/CD\par}

\section{Experience}
\role{Multiplatform Full Stack Software Engineer, XR Lab}{Sep 2026 -- Present}{DICTUC - Pontificia Universidad Católica de Chile (full time)}{Santiago, Chile}
\begin{points}
  \item Moved a three-headset mixed-reality study app to a native \textbf{VIVE Focus Vision} build (\textbf{89.5 fps}, passthrough, hand and eye tracking) for ANID project 13250116 on emotion and collaboration in XR.
  \item Synchronized gaze, motion, and task events with \textbf{MindWare BioLab} physiological recordings to within a few milliseconds, with no turns lost on Wi-Fi drops; built a web panel that runs every headset from a PC or phone.
  \item Built Waves AR (Swift, \textbf{Metal FFT ocean}, ARKit) and a wireless dev loop that cut an AR iteration from \textbf{about 8 min to 4--7 s}.
\end{points}

\role{Freelance Software Engineer, Physical AI \& Robotics}{Sep 2026 -- Present}{TEK Chile (part time)}{Santiago, Chile}
\begin{points}
  \item Built the control app for a \textbf{Dobot CR5S} roll-packing cell (Raspberry Pi, tablet UI, TCP/IP and Modbus) where one plan drives both the 3D preview and the arm, plus an \textbf{Isaac Sim} digital twin and YOLO26 roll counting and box-corner vision.
  \item Built a browser controller with server-side safety for a \textbf{Dobot Rover X1} inspection robot: 18 real patrol runs with no contact, video raised from 5 to about 20 fps, Pi latency cut from \textbf{125 to about 25 ms}, and a 0.48 s emergency stop.
\end{points}

\role{Full Stack Developer}{Sep 2025 -- Sep 2026}{Dily - Automotive Fintech Platform}{Santiago, Chile}
\begin{points}
  \item Built lending workflows (\textbf{TypeScript/Node.js}, Angular/Nx, PostgreSQL, AWS) with DDD, hexagonal architecture, and tests.
\end{points}

\role{Lead Engineer, Alerting Platform Migration}{2026}{GridWorks}{Santiago, Chile}
\begin{points}
  \item Led the staged rebuild of an industrial alerting platform (Next.js, Node.js, MQTT, PostgreSQL) without interrupting live workflows.
\end{points}

\role{Software Engineer}{Dec 2024 -- Jul 2025}{Flair - Commercial-Building Energy Optimization}{Santiago, Chile}
\begin{points}
  \item Python microservices, Vue.js UI, and InfluxDB time-series workflows on AWS, tied to \textbf{HVAC energy reductions of up to 50\%}.
\end{points}

\role{Research Software Engineer}{Jul 2023 -- Jul 2025}{Politiktok (Fondecyt) - Pontificia Universidad Católica de Chile}{Santiago, Chile}
\begin{points}
  \item GPU-backed Python/PostgreSQL pipelines for \textbf{100,000+ records}, about \textbf{10x faster} through parallelization; Springer co-author.
\end{points}

\role{Simulation Engineer}{Jan 2023 -- Jul 2024}{Pontificia Universidad Católica de Chile - Geoscience}{Santiago, Chile}
\begin{points}
  \item Optimized \textbf{CUDA} parallel algorithms for tectonic-plate modeling and seismic-risk analysis, reaching a \textbf{15x speedup}.
\end{points}

\role{AI Systems Developer}{Jan 2024 -- Mar 2024}{Drone Response - University of Notre Dame}{Notre Dame, IN, USA}
\begin{points}
  \item Mission planning for autonomous multi-drone coordination (Python, MQTT, OpenAI API, Java Spring); NASA/NSF-backed.
\end{points}

\section{Teaching \& Recognition}
\entry{Teaching}{Senior TA at PUC, 14+ roles in OS, HPC, Networks, XR (2023--2025); Python instructor, Escuela Militar (2024--2026).}
\entry{Conferences}{SC26 Early Career Program with travel grant (Chicago, 2026); ACM SIGGRAPH 2026 Student Volunteer.}

\section{Education}
\role{Pontificia Universidad Católica de Chile}{Mar 2020 -- Jul 2025}{B.S. Computer Science Engineering, Major in Software Engineering, Minor in Data Science; with Distinction}{Santiago, Chile}

\end{document}
